\section{Rappels sur les règles de calcul}
\subsection{Vocabulaire}
\begin{definitionbox}
    On distingue trois opérations:
    \begin{itemize}[label=\textbullet]

        \item \bb{Addition} : \hl[VertClair]{4} + \hl[VertClair]{5} = \hl[RougeClair]{9}
            \begin{itemize}[label=\textbullet]
                \item \hl[VertClair]{4} et \hl[VertClair]{5} sont
                    appelés les \answer{termes}
                \item \hl[RougeClair]{9} est appelé la somme
            \end{itemize} 


        \item \bb{Soustraction}: \hl[VertClair]{12} - \hl[VertClair]{5}
            = \hl[RougeClair]{7}
            \begin{itemize}[label=\textbullet]
                \item \hl[VertClair]{12} et \hl[VertClair]{5} sont
                    appelés les \answer{termes}
                \item \hl[RougeClair]{7} est appelé la \answer[RougeClair]{différence}
            \end{itemize} 

        \item \bb{Multiplication}: \hl[RougeClair]{12} $\times$ \hl[RougeClair]{4}
            = \hl[VertClair]{7}
            \begin{itemize}[label=\textbullet]
                \item \hl[RougeClair]{12} et \hl[RougeClair]{4} sont
                    appelés les \answer[RougeClair]{facteurs}
                \item \hl[VertClair]{48} est appelé le
                    \answer[VertClair]{produit}
            \end{itemize} 

        \item \bb{Division} : \hl[JauneClair]{18} $\div$
            \hl[RougeClair]{3} = \hl[VertClair]{6}
            \begin{itemize}[label=\textbullet]
                \item \hl[JauneClair]{18} est appelé le
                    \answer[JauneClair]{dividende}.
                \item \hl[RougeClair]{3} est 
                    appelé le \answer[RougeClair]{diviseur}.

                \item \hl[VertClair]{6} est appelé le \answer{quotient}
            \end{itemize} 


    \end{itemize}
\end{definitionbox}

\subsection{Rappel sur les règles de calcul}
\begin{propbox}
    Pour effectuer des opérations \bb{avec et sans} les parenthèses on se doit de
    respecter \bb{6 règles}:
    \begin{itemize}[label=\textbullet]
        \item \bb{Règle 1} : Lorsqu'il n'y a que des additions et des
            soustractions, on effectue les calculs \bb{de la gauche vers la droite}.
        \item \bb{Règle 2} : Lorsqu'il n'ya que des multiplication et des
            divisions, on effectue les calculs \bb{de la gauche vers la droite}.
        \item \bb{Règle 3} : La \bb{multiplication} est \bb{prioritaire} devant
            l'addition et la soustraction.
        \item \bb{Règle 4} : La \bb{division} est \bb{prioritaire} devant
            l'addition et la soustraction.
        \item \bb{Règle 5} : On \bb{commence} par effectuer les \bb{calculs entre
            parenthèses.}
        \item \bb{Règle 6} : On \bb{commence} par effectuer les calcules entre
            \bb{les parenthèses les plus intérieures}.
    \end{itemize}
\end{propbox}

\begin{exemplebox}
    \begin{itemize}[label=\textbullet]
        \item \bb{Règle 1 :} avec uniquement des additions et des soustractions,
        on calcule de la gauche vers la droite :
        \[
            18-5+7
            =\answer{13 +7}
            =\answer{20}.
        \]

        \item \bb{Règle 2 :} avec uniquement des multiplications et des divisions,
        on calcule de la gauche vers la droite :
        \[
            48\div4\times3
            =\answer{12\times 3}
            =\answer{36}.
        \]

        \item \bb{Règle 3 :} la multiplication est prioritaire devant
        l'addition et la soustraction :
        \[
            7+3\times5
            =\answer{7+15}
            =\answer{22}.
        \]

        \item \bb{Règle 4 :} la division est prioritaire devant
        l'addition et la soustraction :
        \[
            20-12\div3
            =\answer{20-4}
            =\answer{16}.
        \]

        \item \bb{Règle 5 :} on commence par effectuer les calculs entre
        parenthèses :
        \[
            (8+4)\times3
            =\answer{12\times3}
            =\answer{36}.
        \]

        \item \bb{Règle 6 :} lorsqu'il y a plusieurs parenthèses imbriquées,
        on commence par les parenthèses les plus intérieures :
        \[
            5\times(3+(8-5))
            =\answer{5\times(3+3)}
            =\answer{5\times6}
            =\answer{30}.
        \]
    \end{itemize}

    \medskip

    \textbf{Un exemple pour tout retenir :}
    \[
        4+2\times(10-(6+2))
        =\answer{4+2\times(10-8)}
        =\answer{4+2\times2}
        =\answer{4+4}
        =\answer{8}.
    \]

\end{exemplebox}

\begin{attentionbox}
    \bb{Attention :} les multiplications et les divisions sont prioritaires
    devant les additions et les soustractions. La multiplication et la division
    ont \bb{la même priorité} : on calcule alors de la gauche vers la droite.

\end{attentionbox}


\begin{propbox}
    Dans un calcul, la dernière opération nous dit s'il s'agit d'une \bb{somme},
    d'une \bb{différence}, d'un \bb{produit} ou d'un \bb{quotient}
\end{propbox}

\begin{exemplebox}
    \begin{itemize}[label=\textbullet]
    \item Dans le calcul  
    \[
        7+3\times4
    \]
    la dernière opération effectuée est une addition :
    \[
        7+3\times4
        =7+\answer{12}
        =\answer{19}.
    \]
    Le résultat final est donc un \answer{somme}.
    
    \item Dans le calcul $$20-(5+3)$$

    on commence par calculer la parenthèse :
    \[
        20-(5+3)
        =20-\answer{8}
        =\answer{12}.
    \]
    La dernière opération est une \answer{soustraction} : le résultat final est donc une
    \answer{différence}.

    \item Dans le calcul $$(6+4)\times5$$
    on commence par calculer la parenthèse :
    \[
        (6+4)\times5
        =\answer{10}\times5
        =\answer{50}.
    \]
    La dernière opération est une multiplication : le résultat final est donc un
    \answer{produit}.

    \item Dans le calcul $$ 24\div(3+1) $$
    on commence par calculer la parenthèse :
    \[
        24\div(3+1)
        =24\div\answer{4}
        =\answer{6}.
    \]
    La dernière opération est une division : le résultat final est donc un
    \answer{quotient}.
\end{itemize}

\end{exemplebox}


